Dans le référentiel terrestre, $A$ est la position initiale ($v_G=v_0$) et $B$
la position d'arrêt ($v_G=0$).

\begin{figure}[H]
    \centering
    \begin{tikzpicture}[]
        \newdimen\el \el=6cm
        \Sol{1.5\el}{5mm}{sol}
        \foreach \x/\p/\t in {-0.5\el/A/$A~(v_{0})$,0/G/,0.5\el/B/$B~(v=0)$}{
            \node[solide={1cm/1cm},label=\t,anchor=south]
                (\p) at ([xshift=\x]sol.north) {};
            \fill (\p.center) circle (2pt);
        }
        \foreach \P [count=\i] in {A.center,B.center}
            \draw[densely dashed] (\P) -- ++(0,-2) coordinate(E\i);
        \draw[<->] (E1) -- node[below] {$\ell$} (E2);
        \foreach \e in {0.1,0.3,...,1}
            \draw[thick] ([xshift=-3pt]$(A.north west)!\e!(A.south west)$) --
                ++(-0.4,0);
        \foreach \e in {0.1,0.3,...,1}
            \draw[thick] ([xshift=-3pt]$(G.north west)!\e!(G.south west)$) --
                ++(-0.2,0);
        \draw[-Latex,very thick,Green] ([xshift=2pt]G.south) -- ++(0,1.8)
            node[below right,inner ysep=0pt] {$\vec{R}_{N}$};
        \draw[-Latex,very thick,red] ([xshift=2pt]G.south) -- ++(-1.8,0)
            node[above right,inner xsep=0pt] {$\vec{f}$};
        \draw[-Latex,very thick,blue] (G.center) -- ++(0,-1.8)
            node[above right,inner ysep=0pt] {$\vec{P}$};
    \end{tikzpicture}
\end{figure}



\begin{itemize}
    \item La variation d'énergie cinétique entre les positions $A$ et $B$ vaut~:
          \[
              \Delta E_c=E_c(B)-E_c(A)=0-\frac{1}{2}m v_0^2=-\frac{1}{2}m v_0^2.
          \]
          
    \item Le travail des forces extérieures appliquées est~:
          \[
              W_{AB}(\vec{P})=W_{AB}(\vec{R}_N)=0~; \quad
              W_{AB}(\vec{f})=f\times AB\times\cos(\ang{180})=-f\ell.
          \]
\end{itemize}


D'où~:
\[
    \Delta E_c=\sum W_{AB}(\vec{F}_{\text{ext}})
    \quad\Rightarrow\
    -\frac{1}{2} m v_0^2=-f \ell
    \quad\Rightarrow\
    \mathbox{f=\frac{m v_0^2}{2 \ell}}
\]

Application numérique~:
\[
    f=\frac{0,920 \times(2,0)^2}{2 \times 0,80}=\qty{2.3}{\N}.
\]
